\subsection*{Общая характеристика работы}

\newcommand{\actuality}{\underline{\textbf{Актуальность темы.}}}
\newcommand{\aim}{\underline{\textbf{Целью}}}
\newcommand{\tasks}{\underline{\textbf{задачи}}}
\newcommand{\defpositions}{\underline{\textbf{Основные положения, выносимые на~защиту:}}}
\newcommand{\novelty}{\underline{\textbf{Научная новизна:}}}
\newcommand{\influence}{\underline{\textbf{Практическая значимость}}}
\newcommand{\reliability}{\underline{\textbf{Достоверность}}}
\newcommand{\probation}{\underline{\textbf{Апробация работы.}}}
\newcommand{\contribution}{\underline{\textbf{Личный вклад.}}}
\newcommand{\publications}{\underline{\textbf{Публикации.}}}

{\actuality}
Данная работа посвящена разработке рентгенооптических трактов синхротронного источника СКИФ --- <<Сибирский кольцевой источник фотонов>>. За последние три десятилетия мир увидел активное развитие специализированных источников синхротронного излучения и соответствующих методов исследования вещества с использованием синхротронного излучения в рентгеновском диапазоне. Главные параметры излучения, который достигаются на данных установках является высокий поток фотонов, направленность излучения в малый телесный угол, когерентность. Эти параметры крайне необходимы для проведения качественных экспериментов с революционными результатами в области химии, биологии, материаловедении, медицины и многих других отраслях науки и техники.

Высокая востребованность данной работы заключается в том, что в отечественной науке наблюдается стагнация в области развития специализированных источников рентгеновского излучения. Проектируемый в Новосибирске синхротронный источник является первым на территории России специализированным источником с проектными параметрами не уступающими мировым установкам, а по некоторым данным с запасом превосходящих их, такие, например, как: MAX-IV, NSLS-II, PETRA-III, Diamond и д.р.

\textbf{Цель} данной работы --- разработка проекта станций первой очереди, вставными устройствами на которых являются сверхпроводящие ондуляторы. Это станции: 1-1 --- <<Микрофокус>>, 1-2 --- <<Структурная диагностика>>, 1-4 --- <<XAFS-спектроскопия и магнитный дихроизм>>.  

В цели проектирования входит ряд \textbf{задач}:
\begin{itemize}

	\item Расчёт ондуляторного излучения с помощью численного моделирования, получение спектров и сечений пучка из указанных устройств, максимально обективно описывающих реальное излучение.
	\item Разработка оптических трактов: расчёт тепловых нагрузок, расчёт спектров и сечений пучка после прохождение оптических элементов. 
	\item Разработка программного кода для реализации выше приведённых задач и удобному воспроизведению результатов расчётов любым участником проекта.

\end{itemize} % Характеристика работы по структуре во введении и в автореферате не отличается (ГОСТ Р 7.0.11, пункты 5.3.1 и 9.2.1), потому её загружаем из одного и того же внешнего файла, предварительно задав форму выделения некоторым параметрам

%Диссертационная работа была выполнена при поддержке грантов ...

%\underline{\textbf{Объем и структура работы.}} Диссертация состоит из~введения, четырех глав, заключения и~приложения. Полный объем диссертации \textbf{ХХХ}~страниц текста с~\textbf{ХХ}~рисунками и~5~таблицами. Список литературы содержит \textbf{ХХX}~наименование.

%\newpage
\subsection*{Содержание работы}
Во \underline{\textbf{введении}} обосновывается актуальность исследований, проводимых в рамках данной диссертационной работы, приводится обзор научной литературы по изучаемой проблеме, формулируется цель, ставятся задачи работы, сформулированы научная новизна и практическая значимость представляемой работы.

\underline{\textbf{Первая глава}} посвящена ...

 картинку можно добавить так:
\begin{figure}[ht] 
  \center
  \includegraphics [scale=0.27] {latex}
  \caption{Подпись к картинке.} 
  \label{img:latex}
\end{figure}

Формулы в строку без номера добавляются так:
\[ 
  \lambda_{T_s} = K_x\frac{d{x}}{d{T_s}}, \qquad
  \lambda_{q_s} = K_x\frac{d{x}}{d{q_s}},
\]

\underline{\textbf{Вторая глава}} посвящена исследованию 

\underline{\textbf{Третья глава}} посвящена исследованию 

В \underline{\textbf{четвертой главе}} приведено описание 

В \underline{\textbf{заключении}} приведены основные результаты работы, которые заключаются в следующем:
%% Согласно ГОСТ Р 7.0.11-2011:
%% 5.3.3 В заключении диссертации излагают итоги выполненного исследования, рекомендации, перспективы дальнейшей разработки темы.
%% 9.2.3 В заключении автореферата диссертации излагают итоги данного исследования, рекомендации и перспективы дальнейшей разработки темы.
\begin{enumerate}
  \item Были выработаны оптимальные критерии отбора событий процесса ${e^+\:e^- \to K_{S}\:K_{L}\:\pi^0}$. Вклад данного процесса в общее адронное сечение менее 5\%, тем не менее удалось набрать статистику в 1924 событий при помощи выбора оптимальных условий отбора.
  \item Определена эффективность регистрации при помощи Монте-Карло моделирования, которая составила около 2\% в диапозоне энергий от \text{1.1 ГэВ} до \text{1.4 ГэВ} и около 4.5\% для энергий от 1.4 ГэВ до 2 ГэВ со статистической ошибкой $\sim1\%$.
  \item Измерено сечение процесса ${e^+\:e^- \to K_{S}\:K_{L}\:\pi^0}$ в диапозоне энергий от \text{1.1 ГэВ} до \text{2 ГэВ} со статистической погрешностью около 10\%. Вычислены радиационные поправки с точностью $\sim1\%$, с последующим получением конечного результата --- борновское сечение изучаемого процесса.
  \item Сделана оценка вклада трех основных фоновых событий, общий вклад которых составил не более 10\% от сечения процесса ${e^+\:e^- \to K_{S}\:K_{L}\:\pi^0}$.
\end{enumerate}



%\newpage
\renewcommand{\refname}{\large Публикации автора по теме диссертации}
%\nocite{*}
%\insertbiblioauthor                          % Подключаем Bib-базы
\insertbibliofull
